\documentclass[11pt]{article}
\usepackage{mathblatt}
\def\mitzone{1}
\begin{document}
\blattkopf{Punkte und Strecken im Koordinatensystem}{Gesamt}
\verzeichniszeile{\verz{zone}{Kennst du schon (Nr.~1–10)} \verztrenn \verz{e1}{1 Punkte darstellen und Lage lesen (Nr.~11–12)} \verztrenn \verz{e2}{2 Streckenlängen, Mittelpunkt und Teilpunkte (Nr.~13–14)} \verztrenn \verz{e3}{3 Dreiecke nachweisen (Nr.~15–16)} \verztrenn \verz{e4}{4 Vierecke nachweisen (Nr.~17–18)} \verztrenn \verz{e5}{5 Körper im Koordinatensystem und Drehungen (Nr.~19–20)} \verztrenn \verz{abhaken}{Das kann ich}}
% Zone „Das kennst du schon“ – aus bank/punkte-und-strecken-im-koordinatensystem/zone.jsonl (zusammenbau v0.1)
\einheitenkopf[zone][Kennst du schon]{Das kennst du schon}
\setcounter{aufgabe}{0}

%% TODO Titel: Ich-kann-Satz fehlt in der Bank (Platzhalter: Kettenname) – Nr. 1
%% TODO Anweisung: ein Satz über der ersten Teilaufgabe fehlt in der Bank – Nr. 1
\begin{aufgabe}{Punkte im ebenen Koordinatensystem mit vier Quadranten lesen und eintragen, Vorzeichen der Koordinaten deuten}
\begin{teile}
\teil Trage $A(3 | 2)$ ein. Lies $B$ ab.

\begin{ksys}[xmin=-5,xmax=5,ymin=-5,ymax=5] \punkt{1}{4}{B} \end{ksys}
B( \leerfeld | \leerfeld )
\teil Trage $E(-3 | 2)$ und $F(1{,}5 | -2)$ ein. Lies $G$ ab.

\begin{ksys}[xmin=-5,xmax=5,ymin=-5,ymax=5] \punkt{-2}{-1}{G} \end{ksys}
G( \leerfeld | \leerfeld )
\end{teile}
\end{aufgabe}

%% TODO Titel: Ich-kann-Satz fehlt in der Bank (Platzhalter: Kettenname) – Nr. 2
%% TODO Anweisung: ein Satz über der ersten Teilaufgabe fehlt in der Bank – Nr. 2
\begin{aufgabe}{Schrägbilder von Quadern, Prismen und Pyramiden lesen und skizzieren, Netze aufklappen}
\begin{teile}
\teil Quader im Bild: Wie viele Kanten, wie viele davon verdeckt?

\quader{4}{2}{3}{a}{b}{c}
Kanten: \leerfeld , verdeckt: \leerfeld
\teil Zeichne ein Schrägbild des Quaders mit $5$ cm Breite, $3$ cm Tiefe und $2$ cm Höhe (Tiefe schräg und halb so lang).

\begin{ksys}[xmin=0,xmax=8,ymin=0,ymax=5] \end{ksys}
\end{teile}
\end{aufgabe}

%% TODO Titel: Ich-kann-Satz fehlt in der Bank (Platzhalter: Kettenname) – Nr. 3
%% TODO Anweisung: ein Satz über der ersten Teilaufgabe fehlt in der Bank – Nr. 3
\begin{aufgabe}{Verbindungsvektor bilden (Spitze minus Fuß), Vektoren addieren und vervielfachen, den Betrag berechnen}
\begin{teile}
\teil $A(2 | 1 | 3)$, $B(5 | 5 | 3)$: $\overrightarrow{AB}$ und seine Länge?
\teil $\vec a = (2 | -1 | 3)$, $\vec b = (-4 | 0 | 1{,}5)$: $2\vec a + \vec b$?
\end{teile}
\end{aufgabe}

%% TODO Titel: Ich-kann-Satz fehlt in der Bank (Platzhalter: Kettenname) – Nr. 4
%% TODO Anweisung: ein Satz über der ersten Teilaufgabe fehlt in der Bank – Nr. 4
\begin{aufgabe}{Satz des Pythagoras anwenden und mit Wurzeln rechnen (Wurzel isolieren, dann quadrieren)}
\begin{teile}
\teil Katheten $6$ cm und $8$ cm: Hypotenuse? \leerfeld[cm]
\teil Hypotenuse $2{,}5$ m, eine Kathete $1{,}5$ m: andere Kathete? \leerfeld[m]
\end{teile}
\end{aufgabe}

%% TODO Titel: Ich-kann-Satz fehlt in der Bank (Platzhalter: Kettenname) – Nr. 5
%% TODO Anweisung: ein Satz über der ersten Teilaufgabe fehlt in der Bank – Nr. 5
\begin{aufgabe}{Flächenformeln (Dreieck, Trapez), Prozentrechnung und Strahlensätze für die Sachanhänge}
\begin{teile}
\teil Dreieck mit Grundseite $6$ cm und Höhe $4$ cm: Flächeninhalt? \leerfeld[cm²]
\teil Trapez mit $a = 5{,}5$ m, $c = 2{,}5$ m, $h = 1{,}2$ m: Flächeninhalt? \leerfeld[m²]
\end{teile}
\end{aufgabe}

%% TODO Titel: Ich-kann-Satz fehlt in der Bank (Platzhalter: Kettenname) – Nr. 6
%% TODO Anweisung: ein Satz über der ersten Teilaufgabe fehlt in der Bank – Nr. 6
\begin{aufgabe}{Skalarprodukt in Koordinaten berechnen und Orthogonalität als „Skalarprodukt null“ erkennen}
\begin{teile}
\teil $(2 | 1 | 3) \circ (1 | 4 | 2)$? \leerfeld
\teil $(2 | -3 | 1) \circ (-1 | 2 | 0{,}5)$? \leerfeld
\end{teile}
\end{aufgabe}

%% TODO Titel: Ich-kann-Satz fehlt in der Bank (Platzhalter: Kettenname) – Nr. 7
%% TODO Anweisung: ein Satz über der ersten Teilaufgabe fehlt in der Bank – Nr. 7
\begin{aufgabe}{Eigenschaften der Dreiecke und Vierecke kennen (gleichschenklig, gleichseitig, rechtwinklig; Parallelogramm, Rechteck, Raute, Quadrat, Trapez, Drachenviereck) und den Satz des Thales}
\begin{teile}
\teil Welches Viereck hat vier gleich lange Seiten und vier rechte Winkel? \\ \kreuz{Raute} \\ \kreuz{Quadrat} \\ \kreuz{Rechteck}
\teil Zwei Paare paralleler Gegenseiten, kein rechter Winkel, verschieden lange Nachbarseiten: welches Viereck? \\ \kreuz{Trapez} \\ \kreuz{Parallelogramm} \\ \kreuz{Drachenviereck}
\end{teile}
\end{aufgabe}

%% TODO Titel: Ich-kann-Satz fehlt in der Bank (Platzhalter: Kettenname) – Nr. 8
%% TODO Anweisung: ein Satz über der ersten Teilaufgabe fehlt in der Bank – Nr. 8
\begin{aufgabe}{Kollinearität zweier Vektoren prüfen (Vielfaches)}
\begin{teile}
\teil Ist $(4 | 2 | 6)$ ein Vielfaches von $(2 | 1 | 3)$? Faktor? Faktor \leerfeld
\teil Ist $(4 | -2 | -1)$ ein Vielfaches von $(-2 | 1 | 0{,}5)$? Faktor? Faktor \leerfeld
\end{teile}
\end{aufgabe}

%% TODO Titel: Ich-kann-Satz fehlt in der Bank (Platzhalter: Kettenname) – Nr. 9
%% TODO Anweisung: ein Satz über der ersten Teilaufgabe fehlt in der Bank – Nr. 9
\begin{aufgabe}{Verbindungsvektor bilden (Spitze minus Fuß), Vektoren addieren und vervielfachen, den Betrag berechnen – Fehler finden}
\begin{teile}
\teil Ich finde den Fehler: $U(1 | 2 | 2)$, $V(3 | 3 | 4)$. Ali rechnet $|\overrightarrow{UV}| = \sqrt{3^2 + 3^2 + 4^2} = \sqrt{34}$.
\end{teile}
\end{aufgabe}

%% TODO Titel: Ich-kann-Satz fehlt in der Bank (Platzhalter: Kettenname) – Nr. 10
%% TODO Anweisung: ein Satz über der ersten Teilaufgabe fehlt in der Bank – Nr. 10
\begin{aufgabe}{Verbindungsvektor bilden (Spitze minus Fuß), Vektoren addieren und vervielfachen, den Betrag berechnen – selbst rechnen}
\begin{teile}
\teil $K(2 | 1 | 5)$, $L(6 | 9 | 6)$: Länge von $KL$? \leerfeld
\end{teile}
\end{aufgabe}

\clearpage
% Einheit 1 von 5 – Punkte darstellen und Lage lesen (bank/punkte-und-strecken-im-koordinatensystem/e1.jsonl, zusammenbau v0.1)
%% TODO Zeitmarke Einheit 1: Marken-Zeile nicht lesbar
%% TODO Zweigzeile Teil 1: was hier gelernt wird (Satz in Schülersprache aus dem Einheitstitel) – Einheit 1
\einheitenkopf[e1]{Einheit 1 von 5 · Punkte darstellen und Lage lesen}
\zweigzeile{Abitur GK}
\setcounter{aufgabe}{10}

%% TODO Titel: Ich-kann-Satz fehlt in der Bank (Platzhalter: Kettenname) – Nr. 11
%% TODO Anweisung: ein Satz über der ersten Teilaufgabe fehlt in der Bank – Nr. 11
\begin{aufgabe}{Darstellen und Lage lesen}
%% TODO \rechenplatz{n}: Schrittzahl des Grundfalls fehlt in der Bank (nur bei fester Schreibform, 2.3 b)
\begin{teile}
\teil $P(4 | -1 | 0)$ – was sagt die Null? \\ \kreuz{$P$ liegt auf der $x_3$-Achse} \\ \kreuz{$P$ liegt in der $x_1x_2$-Ebene} \\ \kreuz{$P$ liegt in der $x_2x_3$-Ebene}
\teil $A(5 | 3 | 1)$ und $B(2 | 1 | 4)$ eintragen, $C$ ablesen.

\begin{ksys3} \rpunkt{4}{5}{2}{C} \end{ksys3}
C( \leerfeld | \leerfeld | \leerfeld )
\teil Ein Holzwürfel mit Kante $4$ hat eine Ecke im Ursprung und liegt im ersten Oktanten (alle Koordinaten $\ge 0$). Zeichne das Viereck $P(4 | 0 | 1)$, $Q(1 | 4 | 0)$, $R(0 | 4 | 2)$, $S(3 | 0 | 4)$ ein. \hfill (Abitur 2019 GK)

\begin{ksys3} \rquader{0,0,0}{4}{4}{4} \end{ksys3}
\teil Zeichne das Viereck $A(5 | -3 | 0)$, $B(5 | 2 | 0)$, $C(-1 | 4 | 4)$, $D(-1 | -1 | 4)$ in das Koordinatensystem. \hfill (Abitur 2018 LK)

\begin{ksys3} \end{ksys3}
\teil Ein Quader hat eine Ecke im Ursprung $O$ und die Kanten $5$, $6$ und $3$ entlang der $x_1$-, $x_2$- und $x_3$-Achse. $P(5 | 2 | 3)$ liegt auf einer Kante. Zeichne das Dreieck mit den Ecken $O$, $P$ und $R(0 | 6 | 3)$ ein. \hfill (Abitur 2026 GK)

\begin{ksys3} \rquader{0,0,0}{5}{6}{3} \end{ksys3}
\teil Begründe: $P(-3 | 4 | 3)$ und $Q(5 | -1 | 2)$ liegen auf derselben Seite der $x_1x_2$-Ebene. \hfill (Abitur 2025 GK)
\teil Der Grundriss eines Pavillons ist ein Achteck in der $x_1x_2$-Ebene, symmetrisch zur $x_1$- und zur $x_2$-Achse. Gegeben sind $A(6 | 0)$, $B(5 | 3)$, $C(0 | 4)$. Zeichne das vollständige Achteck. \hfill (Abitur 2025 GK)

\begin{ksys}[xmin=-7,xmax=7,ymin=-7,ymax=7,xlabel=$x_1$,ylabel=$x_2$] \punkt{6}{0}{A} \punkt{5}{3}{B} \punkt{0}{4}{C} \end{ksys}
\teil Eine Pyramide hat den quadratischen Boden $ABCD$ mit der Seite $3$ cm; die Spitze $S$ liegt $4$ cm senkrecht über $A$. Im Gitter sind der Boden und die Spitze des Seitendreiecks $ABS$ markiert. Zeichne die drei übrigen Seitendreiecke ein. \hfill (Abitur 2024 GK)

\begin{ksys}[xmin=-5,xmax=9,ymin=-5,ymax=9] \punkt{0}{0}{A} \punkt{3}{0}{B} \punkt{3}{3}{C} \punkt{0}{3}{D} \punkt{0}{-4}{S} \end{ksys}
\teil $B(5 | 1 | 4)$ und $C(1 | 4 | 3)$ sind Ecken des Parallelogramms $ABCD$, $M(2 | 2 | 2)$ ist der Schnittpunkt seiner Diagonalen. Alle Punkte werden parallel zur $x_3$-Achse in die $x_1x_2$-Ebene verschoben. Zeichne $A'B'C'D'$ und $M'$ ein. \hfill (Abitur 2024 LK)

\begin{ksys3} \end{ksys3}
\end{teile}
\end{aufgabe}

%% TODO Titel: Ich-kann-Satz fehlt in der Bank (Platzhalter: Kettenname) – Nr. 12
%% TODO Anweisung: ein Satz über der ersten Teilaufgabe fehlt in der Bank – Nr. 12
\begin{aufgabe}{Darstellen und Lage lesen – Fehler finden, Begründen}
\begin{teile}
\teil Ich finde den Fehler: Lena trägt $P(1 | -4 | 2)$ ein: $1$ in $x_1$-Richtung, $4$ in $x_2$-Richtung, $2$ nach oben.
\teil Begründe: $P(4 | 0 | 3)$ liegt in der $x_1x_3$-Ebene.
\end{teile}
\end{aufgabe}

\clearpage
% Einheit 2 von 5 – Streckenlängen, Mittelpunkt und Teilpunkte (bank/punkte-und-strecken-im-koordinatensystem/e2.jsonl, zusammenbau v0.1)
%% TODO Zeitmarke Einheit 2: Marken-Zeile nicht lesbar
%% TODO Zweigzeile Teil 1: was hier gelernt wird (Satz in Schülersprache aus dem Einheitstitel) – Einheit 2
\einheitenkopf[e2]{Einheit 2 von 5 · Streckenlängen, Mittelpunkt und Teilpunkte}
\zweigzeile{Abitur GK}
\setcounter{aufgabe}{12}

%% TODO Titel: Ich-kann-Satz fehlt in der Bank (Platzhalter: Kettenname) – Nr. 13
%% TODO Anweisung: ein Satz über der ersten Teilaufgabe fehlt in der Bank – Nr. 13
\begin{aufgabe}{Streckenlänge, Mittelpunkt, Teilpunkte}
%% TODO \rechenplatz{n}: Schrittzahl des Grundfalls fehlt in der Bank (nur bei fester Schreibform, 2.3 b)
\begin{teile}
\teil Wo muss der Haken an der Stange zwischen $A$ und $B$ hängen, damit er genau in der Mitte ist? Was ist gefragt? \\ \kreuz{eine Länge} \\ \kreuz{ein Punkt} \\ \kreuz{eine Figur}
\teil $A(2 | 1 | 3)$, $B(4 | 4 | 9)$: Länge von $AB$? \leerfeld LE
\teil Ein Zeltmast endet in $S(2 | 2 | 6)$; vier gleich lange Leinen führen zu Heringen am Boden, einer davon steckt in $H(4 | 5 | 0)$ (1 LE = 1 m). Je Leine kommen $40$ cm für die Knoten dazu. Wie viel Leine braucht man insgesamt? \leerfeld[m]
\teil Ein Gummiseil ist von $A(3 | 1 | 2)$ nach $B(9 | 9 | 2)$ gespannt (1 LE = 1 m) und dabei um $25\,\%$ länger als ungespannt. Wie lang ist es ungespannt? \hfill (Abitur 2020 GK) \\ \leerfeld[m]
\teil Eine Pyramide ist symmetrisch zur $x_1x_3$-Ebene; $D(1 | 3 | 0)$ ist eine Grundecke, $S(3 | 0 | 6)$ die Spitze. Zeige: $W(2 | 1{,}5 | 3)$ ist der Mittelpunkt der Kante $DS$. Gib den Mittelpunkt $T$ der zu $DS$ symmetrischen Kante an. \hfill (Abitur 2026 GK) \\ T( \leerfeld | \leerfeld | \leerfeld )
\teil Eine Strebe führt von $P(9 | 0 | 3)$ nach $Q(0 | 3 | 0)$; zwei Schellen teilen sie in drei gleich lange Stücke. Koordinaten der Schellen? \hfill (Abitur 2024 GK)
\teil Eine Stütze führt von $M(3 | 3 | 0)$ zur Spitze $D(0 | 0 | 9)$. $Q_h$ ist der Punkt der Strecke in der Höhe $h$. Koordinaten von $Q_h$? Gib $Q_{3}$ an. \hfill (Abitur 2022 LK)
\teil $A(3 | 5 | 2)$ und $B(5 | 6 | 4)$ liegen auf einer Geraden $g$. Gib einen Punkt $D \ne B$ auf $g$ an, der von $A$ so weit entfernt ist wie $B$. \hfill (Abitur 2020 GK) \\ D( \leerfeld | \leerfeld | \leerfeld )
\teil $A(2 | -1 | 2)$; der Punkt $C$ ist festgelegt durch $\overrightarrow{OC} = 3 \cdot \overrightarrow{OA}$. Länge der Strecke $AC$? \hfill (Abitur 2018 GK) \\ \leerfeld
\teil Die Ecken eines Oktaeders sind die Mittelpunkte der Seitenflächen eines Würfels; $A(2 | 1 | 3)$ und $C(-2 | -3 | 5)$ sind gegenüberliegende Ecken des Oktaeders. Zeige: Die Würfelkante ist $6$ lang. \hfill (Abitur 2024 LK)
\end{teile}
\end{aufgabe}

%% TODO Titel: Ich-kann-Satz fehlt in der Bank (Platzhalter: Kettenname) – Nr. 14
%% TODO Anweisung: ein Satz über der ersten Teilaufgabe fehlt in der Bank – Nr. 14
\begin{aufgabe}{Streckenlänge, Mittelpunkt, Teilpunkte – Fehler finden, Begründen}
\begin{teile}
\teil Ich finde den Fehler: $A(4 | 1 | 3)$, $B(6 | 2 | 5)$. Jan rechnet $|\overrightarrow{AB}| = \sqrt{6^2 + 2^2 + 5^2} = \sqrt{65}$.
\teil Begründe: Die Länge von $\vec v = (a | b | c)$ ist $\sqrt{a^2 + b^2 + c^2}$ – das ist der Satz des Pythagoras im Raum.
\end{teile}
\end{aufgabe}

\clearpage
% Einheit 3 von 5 – Dreiecke nachweisen (bank/punkte-und-strecken-im-koordinatensystem/e3.jsonl, zusammenbau v0.1)
%% TODO Zeitmarke Einheit 3: Marken-Zeile nicht lesbar
%% TODO Zweigzeile Teil 1: was hier gelernt wird (Satz in Schülersprache aus dem Einheitstitel) – Einheit 3
\einheitenkopf[e3]{Einheit 3 von 5 · Dreiecke nachweisen}
\zweigzeile{Abitur GK}
\setcounter{aufgabe}{14}

%% TODO Titel: Ich-kann-Satz fehlt in der Bank (Platzhalter: Kettenname) – Nr. 15
%% TODO Anweisung: ein Satz über der ersten Teilaufgabe fehlt in der Bank – Nr. 15
\begin{aufgabe}{Dreiecke nachweisen}
%% TODO \rechenplatz{n}: Schrittzahl des Grundfalls fehlt in der Bank (nur bei fester Schreibform, 2.3 b)
\begin{teile}
\teil Parallelogramm $ABCD$: Welcher Vektor ist gleich $\overrightarrow{AB}$? \\ \kreuz{$\overrightarrow{CD}$} \\ \kreuz{$\overrightarrow{DC}$} \\ \kreuz{$\overrightarrow{BC}$}
\teil Zeige: Das Dreieck $A(2 | 0 | 1)$, $B(6 | 2 | 1)$, $C(4 | 1 | 5)$ ist gleichschenklig. Nenne Basis und Schenkel. \hfill (Abitur 2023 GK)
\teil Prüfe, ob das Dreieck $A(4 | 2 | 2)$, $B(2 | 4 | 2)$, $C(2 | 2 | 4)$ gleichseitig ist. \hfill (Abitur 2023 GK)
\teil $A(1 | 2 | 0)$, $B(5 | 6 | 0)$ und $C(3 | 4 | p)$ mit $p > 0$. Zeige: Das Dreieck $ABC$ ist für jedes $p$ gleichschenklig mit der Basis $AB$. \hfill (Abitur 2022 GK)
\teil Pyramiden mit $A(0 | 0 | 0)$, $B(7 | 0 | 0)$, $C(0 | 7 | 0)$ und der Spitze $D_p(0 | 0 | p)$, $p > 0$. Begründe ohne Rechnung: Das Dreieck $BCD_p$ ist gleichschenklig. \hfill (Abitur 2022 LK)
\teil Grundfläche $A(6 | 1 | 0)$, $B(1 | 4 | 0)$, $C(1 | -2 | 0)$; Deckkante $E(1 | 3 | 2)$, $F(1 | -1 | 2)$. Zeige: $ABC$ ist gleichschenklig. Begründe: $EF$ ist parallel zur Grundfläche. \hfill (Abitur 2023 GK)
\teil $D(8 | 0 | 2)$, $E(0 | 6 | 2)$, $F(0 | 0 | 2)$ und $M(4 | 3 | 2)$. Begründe: $D$, $E$ und $F$ liegen auf einem Kreis mit dem Mittelpunkt $M$. \hfill (Abitur 2024 GK)
\teil In der $x_1x_2$-Ebene liegen $A(0 | 0)$, $B(6 | 0)$ und $C(3 | 9)$ auf einem Kreis. Koordinaten seines Mittelpunkts $M$? \hfill (Abitur 2019 GK) \\ M( \leerfeld | \leerfeld )
\teil $A(0 | 6 | 1)$, $B(0 | -6 | 1)$ und $C(0 | 0 | p)$ mit $p > 1$. Der Umfang des Dreiecks $ABC$ ist $32$. $p$? \hfill (Abitur 2024 LK) \\ p = \leerfeld
\teil Dreieck $OPQ$ mit dem Ursprung $O$, $P(4 | 2 | 1)$ und $Q(1 | 5 | 3)$. Gib einen Punkt $T \ne P$ an, sodass das Dreieck $OTQ$ denselben Flächeninhalt und denselben Umfang hat. \hfill (Abitur 2023 LK) \\ T( \leerfeld | \leerfeld | \leerfeld )
\teil Dreieck $A(2 | 1 | 2)$, $B(4 | 5 | 4)$, $C(5 | 3 | 1)$ mit $|CA| = |CB|$. Gesucht ist ein Punkt $D \ne C$ mit $|DA| = |DB|$ und $|DA| = |CA|$. Koordinaten von $D$? \hfill (Abitur 2023 GK) \\ D( \leerfeld | \leerfeld | \leerfeld )
\end{teile}
\end{aufgabe}

%% TODO Titel: Ich-kann-Satz fehlt in der Bank (Platzhalter: Kettenname) – Nr. 16
%% TODO Anweisung: ein Satz über der ersten Teilaufgabe fehlt in der Bank – Nr. 16
\begin{aufgabe}{Dreiecke nachweisen – Fehler finden, Begründen}
\begin{teile}
\teil Ich finde den Fehler: Zu zeigen ist, dass $A(2 | 3 | 1)$, $B(6 | 3 | 1)$, $C(4 | 5 | 2)$ ein gleichschenkliges Dreieck mit der Basis $AB$ bilden. Tom: „$|AB| = 4$, $|AC| = 3$ – nicht gleichschenklig.“
\teil Begründe: Für „gleichschenklig“ genügt es, zwei Seitenlängen zu vergleichen. Welche?
\end{teile}
\end{aufgabe}

\clearpage
% Einheit 4 von 5 – Vierecke nachweisen (bank/punkte-und-strecken-im-koordinatensystem/e4.jsonl, zusammenbau v0.1)
%% TODO Zeitmarke Einheit 4: Marken-Zeile nicht lesbar
%% TODO Zweigzeile Teil 1: was hier gelernt wird (Satz in Schülersprache aus dem Einheitstitel) – Einheit 4
\einheitenkopf[e4]{Einheit 4 von 5 · Vierecke nachweisen}
\zweigzeile{Abitur GK}
\setcounter{aufgabe}{16}

%% TODO Titel: Ich-kann-Satz fehlt in der Bank (Platzhalter: Kettenname) – Nr. 17
%% TODO Anweisung: ein Satz über der ersten Teilaufgabe fehlt in der Bank – Nr. 17
\begin{aufgabe}{Vierecke nachweisen}
%% TODO \rechenplatz{n}: Schrittzahl des Grundfalls fehlt in der Bank (nur bei fester Schreibform, 2.3 b)
\begin{teile}
\teil Nachzuweisen: $ABCD$ ist ein Rechteck. Was gehört zum Steckbrief? \\ \kreuz{vier gleich lange Seiten} \\ \kreuz{Parallelogramm und ein rechter Winkel} \\ \kreuz{ein Paar paralleler Seiten}
\teil Zeige: $A(1 | 2 | 0)$, $B(4 | 3 | 2)$, $C(3 | 6 | 3)$, $D(0 | 5 | 1)$ bilden ein Parallelogramm $ABCD$. \hfill (Abitur 2021 LK)
\teil Zeige: $A(2 | 1 | 0)$, $B(4 | 3 | 1)$, $C(2 | 4 | 3)$, $D(0 | 2 | 2)$ bilden ein Rechteck $ABCD$. \hfill (Abitur 2023 GK)
\teil Zeige: $A(0 | 1 | 3)$, $B(2 | 3 | 4)$, $C(3 | 5 | 6)$, $D(1 | 3 | 5)$ bilden eine Raute $ABCD$. \hfill (Abitur 2022 GK)
\teil Zeige: $A(2 | 2 | 2)$, $B(5 | 3 | 2)$, $C(6 | 5 | 4)$, $D(3 | 4 | 4)$ bilden ein Parallelogramm, aber kein Rechteck. \hfill (Abitur 2022 GK)
\teil Eine Kletterwand hat die Ecken $A(4 | 1 | 3)$, $B(1 | 4 | 3)$, $F(1 | 7 | 0)$ und $E(7 | 1 | 0)$ (1 LE = 1 m). Zeige: $ABFE$ ist ein Trapez, in dem zwei gegenüberliegende Seiten gleich lang sind. \hfill (Abitur 2018 GK)
\teil Eine Terrasse ist das Viereck $A(1 | -3 | 0)$, $B(7 | -1 | 0)$, $C(7 | 1 | 0)$, $D(1 | 3 | 0)$ (1 LE = 1 m). Zeige: $ABCD$ ist ein Trapez. Flächeninhalt? \hfill (Abitur 2026 GK) \\ \leerfeld[m²]
\teil Zeige: Das Viereck $K(0 | 0 | 2)$, $L(2 | 3 | 2)$, $M(6 | 0 | 2)$, $N(2 | -3 | 2)$ ist ein Drachenviereck. \hfill (Abitur 2023 GK)
\teil Eine ebene Rasenfläche hat die Ecken $A(2 | 1 | 0)$, $B(9 | 1 | 0)$, $C(8 | 4 | 0{,}5)$, $D(8 | 7 | 1)$ und $E(2 | 7 | 1)$ (1 LE = 1 m). Zeige: $AE$ ist parallel zu $CD$, und $CD$ und $DE$ bilden einen rechten Winkel. \hfill (Abitur 2021 GK)
\teil $P(2 | 1 | 0)$ und $Q(2 | 1 | 5)$ sind benachbarte Ecken eines Quadrats in der Ebene $3x_1 + 4x_2 = 10$. Gib eine weitere Ecke an. \hfill (Abitur 2022 GK)
\teil Ein Quadrat liegt in der $x_2x_3$-Ebene; $P(0 | 6 | 3)$ ist eine Ecke. Der Schnittpunkt seiner Diagonalen liegt auf der Geraden durch $(4 | 3 | 2)$ parallel zur $x_1$-Achse. Zeige: $Q(0 | 2 | 5)$ ist eine der beiden zu $P$ benachbarten Ecken. \hfill (Abitur 2018 LK)
\end{teile}
\end{aufgabe}

%% TODO Titel: Ich-kann-Satz fehlt in der Bank (Platzhalter: Kettenname) – Nr. 18
%% TODO Anweisung: ein Satz über der ersten Teilaufgabe fehlt in der Bank – Nr. 18
\begin{aufgabe}{Vierecke nachweisen – Fehler finden, Begründen}
\begin{teile}
\teil Ich finde den Fehler: $A(1 | 3 | 2)$, $B(4 | 4 | 4)$, $C(5 | 7 | 7)$, $D(2 | 6 | 5)$. Mara: „$\overrightarrow{AB} = (3 | 1 | 2)$, $\overrightarrow{CD} = (-3 | -1 | -2)$ – verschieden, also kein Parallelogramm.“
\teil Begründe: Ist $\overrightarrow{AB} \circ \overrightarrow{AD} \ne 0$, so ist das Parallelogramm $ABCD$ kein Rechteck – mehr muss man nicht prüfen.
\end{teile}
\end{aufgabe}

\clearpage
% Einheit 5 von 5 – Körper im Koordinatensystem und Drehungen (bank/punkte-und-strecken-im-koordinatensystem/e5.jsonl, zusammenbau v0.1)
%% TODO Zeitmarke Einheit 5: Marken-Zeile nicht lesbar
%% TODO Zweigzeile Teil 1: was hier gelernt wird (Satz in Schülersprache aus dem Einheitstitel) – Einheit 5
\einheitenkopf[e5]{Einheit 5 von 5 · Körper im Koordinatensystem und Drehungen}
\zweigzeile{Abitur GK}
\setcounter{aufgabe}{18}

%% TODO Titel: Ich-kann-Satz fehlt in der Bank (Platzhalter: Kettenname) – Nr. 19
%% TODO Anweisung: ein Satz über der ersten Teilaufgabe fehlt in der Bank – Nr. 19
\begin{aufgabe}{Körper und Drehungen}
%% TODO \rechenplatz{n}: Schrittzahl des Grundfalls fehlt in der Bank (nur bei fester Schreibform, 2.3 b)
\begin{teile}
\teil Ein Quader steht auf der $x_1x_2$-Ebene; seine Deckenecke $G$ hat $x_3 = 4$. Was sagt das über die Deckfläche? \\ \kreuz{Sie liegt in der $x_1x_2$-Ebene.} \\ \kreuz{Sie liegt parallel zur $x_1x_2$-Ebene in der Höhe $4$.} \\ \kreuz{Sie steht senkrecht auf der $x_1x_2$-Ebene.}
\teil Gerades Prisma $ABCDEF$ ($D$ über $A$, $E$ über $B$, $F$ über $C$) mit $A(2 | 1 | 0)$, $B(7 | 1 | 0)$, $C(2 | 4 | 0)$, $D(2 | 1 | 3)$: $F$? \hfill (Abitur 2026 GK) \\ F( \leerfeld | \leerfeld | \leerfeld )
\teil Eine Pyramide hat als Grundfläche ein rechtwinkliges Dreieck mit den Katheten $6$ cm und $8$ cm; die Spitze liegt $5$ cm senkrecht über der Ecke mit dem rechten Winkel. Wähle geeignete Koordinaten der vier Ecken (1 LE = 1 cm). \hfill (Abitur 2020 GK)
\teil Ein Würfel mit Kante $2$ hat eine Ecke im Ursprung, alle Ecken haben Koordinaten $\ge 0$. Er wird um $(-1 | -1 | 1)$ verschoben. Welche Ecke hat danach $x_1 < 0$, $x_2 > 0$ und $x_3 > 2$? \hfill (Abitur 2024 GK) \\ ( \leerfeld | \leerfeld | \leerfeld )
\teil Eine Pyramide hat die Grundfläche $A(4 | -1 | 2)$, $B(1 | 3 | 2)$, $C(-1 | 0 | 2)$ und die Spitze $S(2 | 1 | 7)$. Begründe: Die Grundfläche ist parallel zur $x_1x_2$-Ebene. Höhe der Pyramide? \hfill (Abitur 2018 GK) \\ h = \leerfeld
\teil Ein gerader Zylinder mit Radius $5$ und Höhe $6$ steht auf der $x_1x_2$-Ebene; $M(4 | 3 | 6)$ ist der Mittelpunkt der Deckfläche. Prüfe, ob $P(7 | 7 | 0)$ auf dem Rand der Grundfläche liegt. \hfill (Abitur 2020 LK)
\teil Ein Dachgeschoss hat einen rechteckigen Boden von $10$ m mal $8$ m. Die Dachflächen steigen von den beiden Traufwänden (Höhe $0{,}5$ m) gleichmäßig zum First über der Mitte (Höhe $5{,}5$ m). Flächen mit weniger als $2$ m Raumhöhe zählen nicht zur Wohnfläche. Welcher Anteil des Bodens zählt nicht? \hfill (Abitur 2019 GK) \\ \leerfeld \%
\teil Der Punkt $B(0 | 10 | 1)$ wird um die $x_3$-Achse gedreht. Gib drei Bildpunkte an, bei denen eine Koordinate $6$ ist. \hfill (Abitur 2024 LK)
\teil Das Dreieck $A(5 | 3 | 0)$, $B(3 | 5 | 0)$, $C(3 | 3 | 4)$ wird um die Kante $AB$ gedreht, bis $C$ in der $x_1x_2$-Ebene liegt, mit $x_1 < 4$. Koordinaten des gedrehten Punktes $C'$? \hfill (Abitur 2023 GK) \\ C'( \leerfeld | \leerfeld | \leerfeld )
\end{teile}
\end{aufgabe}

%% TODO Titel: Ich-kann-Satz fehlt in der Bank (Platzhalter: Kettenname) – Nr. 20
%% TODO Anweisung: ein Satz über der ersten Teilaufgabe fehlt in der Bank – Nr. 20
\begin{aufgabe}{Körper und Drehungen – Fehler finden, Begründen}
\begin{teile}
\teil Ich finde den Fehler: Gerades Prisma $ABCDEF$ ($D$ über $A$, $E$ über $B$, $F$ über $C$) mit $A(1 | 2 | 0)$, $B(6 | 2 | 0)$, $C(1 | 5 | 0)$ und der Höhe $4$. Paul: $F(6 | 2 | 4)$.
\teil Begründe: Beim geraden Prisma auf der $x_1x_2$-Ebene hat jede Deckenecke dieselbe erste und zweite Koordinate wie ihre Grundecke.
\end{teile}
\end{aufgabe}

% Abhakseite „Das kann ich“ – Zone nur im Gesamt (\mitzone)
%% TODO Ich-kann-Sätze: Platzhalter wie in den Aufgabendateien
\begin{abhakseite}
\ifdefined\mitzone
\abhakgruppe{Kennst du schon}
\abhak{1}{Punkte im ebenen Koordinatensystem mit vier Quadranten lesen und eintragen, Vorzeichen der Koordinaten deuten}
\abhak{2}{Schrägbilder von Quadern, Prismen und Pyramiden lesen und skizzieren, Netze aufklappen}
\abhak{3}{Verbindungsvektor bilden (Spitze minus Fuß), Vektoren addieren und vervielfachen, den Betrag berechnen}
\abhak{4}{Satz des Pythagoras anwenden und mit Wurzeln rechnen (Wurzel isolieren, dann quadrieren)}
\abhak{5}{Flächenformeln (Dreieck, Trapez), Prozentrechnung und Strahlensätze für die Sachanhänge}
\abhak{6}{Skalarprodukt in Koordinaten berechnen und Orthogonalität als „Skalarprodukt null“ erkennen}
\abhak{7}{Eigenschaften der Dreiecke und Vierecke kennen (gleichschenklig, gleichseitig, rechtwinklig; Parallelogramm, Rechteck, Raute, Quadrat, Trapez, Drachenviereck) und den Satz des Thales}
\abhak{8}{Kollinearität zweier Vektoren prüfen (Vielfaches)}
\abhak{9}{Verbindungsvektor bilden (Spitze minus Fuß), Vektoren addieren und vervielfachen, den Betrag berechnen – Fehler finden}
\abhak{10}{Verbindungsvektor bilden (Spitze minus Fuß), Vektoren addieren und vervielfachen, den Betrag berechnen – selbst rechnen}
\fi
\abhakgruppe{Einheit 1 · Punkte darstellen und Lage lesen}
\abhak{11}{Darstellen und Lage lesen}
\abhak{12}{Darstellen und Lage lesen – Fehler finden, Begründen}
\abhakgruppe{Einheit 2 · Streckenlängen, Mittelpunkt und Teilpunkte}
\abhak{13}{Streckenlänge, Mittelpunkt, Teilpunkte}
\abhak{14}{Streckenlänge, Mittelpunkt, Teilpunkte – Fehler finden, Begründen}
\abhakgruppe{Einheit 3 · Dreiecke nachweisen}
\abhak{15}{Dreiecke nachweisen}
\abhak{16}{Dreiecke nachweisen – Fehler finden, Begründen}
\abhakgruppe{Einheit 4 · Vierecke nachweisen}
\abhak{17}{Vierecke nachweisen}
\abhak{18}{Vierecke nachweisen – Fehler finden, Begründen}
\abhakgruppe{Einheit 5 · Körper im Koordinatensystem und Drehungen}
\abhak{19}{Körper und Drehungen}
\abhak{20}{Körper und Drehungen – Fehler finden, Begründen}
\end{abhakseite}

\end{document}
